% === Mistake Annotation Preamble ===
% Include this in annotated student solutions via % === Mistake Annotation Preamble ===
% Include this in annotated student solutions via % === Mistake Annotation Preamble ===
% Include this in annotated student solutions via % === Mistake Annotation Preamble ===
% Include this in annotated student solutions via \input{mistake_preamble}

\usepackage{xcolor}
\usepackage{tcolorbox}
\usepackage{marginnote}
\usepackage{xparse}

% Severity color definitions
\definecolor{conceptual}{RGB}{204, 0, 0}       % red — fundamental misunderstanding
\definecolor{procedural}{RGB}{204, 102, 0}      % orange — wrong method/application
\definecolor{mechanical}{RGB}{204, 180, 0}      % yellow-ish — arithmetic/algebra slip
\definecolor{notational}{RGB}{100, 100, 180}    % muted blue — writing/notation issue

% \mistake{content}{explanation}{tag}{severity}
%
% content:      the erroneous LaTeX expression or text
% explanation:  one-sentence description of the mistake
% tag:          tag from the tag bank (e.g., sign-error, wrong-theorem)
% severity:     one of: conceptual, procedural, mechanical, notational
%
% Renders as: colored underline with a margin note.
% Safe inside math environments when content is math.

\NewDocumentCommand{\mistake}{m m m m}{%
    \begingroup
    \color{#4}%
    \underbrace{#1}_{\text{\tiny #3}}%
    \endgroup
    \marginnote{\textcolor{#4}{\scriptsize \textbf{[#4]} #2}}%
}

% Inline variant for use outside math mode
\NewDocumentCommand{\mistaketext}{m m m m}{%
    \textcolor{#4}{\underline{#1}}%
    \marginnote{\textcolor{#4}{\scriptsize \textbf{[#4]} #2}}%
}

\usepackage{xcolor}
\usepackage{tcolorbox}
\usepackage{marginnote}
\usepackage{xparse}

% Severity color definitions
\definecolor{conceptual}{RGB}{204, 0, 0}       % red — fundamental misunderstanding
\definecolor{procedural}{RGB}{204, 102, 0}      % orange — wrong method/application
\definecolor{mechanical}{RGB}{204, 180, 0}      % yellow-ish — arithmetic/algebra slip
\definecolor{notational}{RGB}{100, 100, 180}    % muted blue — writing/notation issue

% \mistake{content}{explanation}{tag}{severity}
%
% content:      the erroneous LaTeX expression or text
% explanation:  one-sentence description of the mistake
% tag:          tag from the tag bank (e.g., sign-error, wrong-theorem)
% severity:     one of: conceptual, procedural, mechanical, notational
%
% Renders as: colored underline with a margin note.
% Safe inside math environments when content is math.

\NewDocumentCommand{\mistake}{m m m m}{%
    \begingroup
    \color{#4}%
    \underbrace{#1}_{\text{\tiny #3}}%
    \endgroup
    \marginnote{\textcolor{#4}{\scriptsize \textbf{[#4]} #2}}%
}

% Inline variant for use outside math mode
\NewDocumentCommand{\mistaketext}{m m m m}{%
    \textcolor{#4}{\underline{#1}}%
    \marginnote{\textcolor{#4}{\scriptsize \textbf{[#4]} #2}}%
}

\usepackage{xcolor}
\usepackage{tcolorbox}
\usepackage{marginnote}
\usepackage{xparse}

% Severity color definitions
\definecolor{conceptual}{RGB}{204, 0, 0}       % red — fundamental misunderstanding
\definecolor{procedural}{RGB}{204, 102, 0}      % orange — wrong method/application
\definecolor{mechanical}{RGB}{204, 180, 0}      % yellow-ish — arithmetic/algebra slip
\definecolor{notational}{RGB}{100, 100, 180}    % muted blue — writing/notation issue

% \mistake{content}{explanation}{tag}{severity}
%
% content:      the erroneous LaTeX expression or text
% explanation:  one-sentence description of the mistake
% tag:          tag from the tag bank (e.g., sign-error, wrong-theorem)
% severity:     one of: conceptual, procedural, mechanical, notational
%
% Renders as: colored underline with a margin note.
% Safe inside math environments when content is math.

\NewDocumentCommand{\mistake}{m m m m}{%
    \begingroup
    \color{#4}%
    \underbrace{#1}_{\text{\tiny #3}}%
    \endgroup
    \marginnote{\textcolor{#4}{\scriptsize \textbf{[#4]} #2}}%
}

% Inline variant for use outside math mode
\NewDocumentCommand{\mistaketext}{m m m m}{%
    \textcolor{#4}{\underline{#1}}%
    \marginnote{\textcolor{#4}{\scriptsize \textbf{[#4]} #2}}%
}

\usepackage{xcolor}
\usepackage{tcolorbox}
\usepackage{marginnote}
\usepackage{xparse}

% Severity color definitions
\definecolor{conceptual}{RGB}{204, 0, 0}       % red — fundamental misunderstanding
\definecolor{procedural}{RGB}{204, 102, 0}      % orange — wrong method/application
\definecolor{mechanical}{RGB}{204, 180, 0}      % yellow-ish — arithmetic/algebra slip
\definecolor{notational}{RGB}{100, 100, 180}    % muted blue — writing/notation issue

% \mistake{content}{explanation}{tag}{severity}
%
% content:      the erroneous LaTeX expression or text
% explanation:  one-sentence description of the mistake
% tag:          tag from the tag bank (e.g., sign-error, wrong-theorem)
% severity:     one of: conceptual, procedural, mechanical, notational
%
% Renders as: colored underline with a margin note.
% Safe inside math environments when content is math.

\NewDocumentCommand{\mistake}{m m m m}{%
    \begingroup
    \color{#4}%
    \underbrace{#1}_{\text{\tiny #3}}%
    \endgroup
    \marginnote{\textcolor{#4}{\scriptsize \textbf{[#4]} #2}}%
}

% Inline variant for use outside math mode
\NewDocumentCommand{\mistaketext}{m m m m}{%
    \textcolor{#4}{\underline{#1}}%
    \marginnote{\textcolor{#4}{\scriptsize \textbf{[#4]} #2}}%
}